\section{Điều nền tảng này chưa giải quyết}
\label{sec:ch02-dieu-con-lai}

Bốn mục trên cùng dựng một đối tượng: kiến trúc, dữ liệu, điều kiện sử dụng và
văn bản trung gian đều là những chỗ mà một \tn{giải thích hậu nghiệm}{post-hoc explanation}
có thể bám vào. Điều không mục nào cho ta là bằng chứng rằng nó đã bám đúng.

Khoảng cách còn lại có dạng rất cụ thể. Ta có thể che một token, đổi một ví
dụ few-shot hoặc so sánh hai câu trả lời, nhưng mỗi thao tác đều đưa vào một
can thiệp và một giả định về dữ liệu thay thế. Do đó, kết quả chỉ đáng tin
trong phạm vi mà can thiệp ấy còn đại diện cho điều ta muốn hỏi về mô hình.
Đây là lý do Phần~II bắt đầu bằng các phương pháp tạo lời giải thích rồi mới
hỏi cách kiểm tra chúng ở Phần~IV.

Chương~\ref{ch:lime} sẽ mở đầu phần phương pháp bằng LIME, một \tn{mô hình thay
thế cục bộ}{local surrogate} được khớp từ các đầu vào đã nhiễu. Khi dùng LIME cho một LLM, ta sẽ
gặp lại toàn bộ các lựa chọn của chương này: đơn vị nào được coi là đặc trưng,
cách nhiễu \tn{lời nhắc}{prompt} nào còn có nghĩa, và thay đổi đầu ra nào thực sự trả lời
câu hỏi đang đặt ra.
